\scanpage{13}
décrire ainsi. Pour $f\in F_M$, \deleted{soit} on note que \unsure{$'$} si $f=p(\widetilde f)$, $\widetilde f\in\widetilde F_M$, alors \deleted{$\alpha_{\widetilde f}(\underline{\omega}_{\widetilde f})\subset(\det M)^*$ est indépendant du choix de $f$, soit} \inserted{l'application composée}
\diagram{\begin{tikzcd}[ampersand replacement=\&,column sep=large]
\underline{\omega}_f \& \underline{\omega}_{\widetilde f}
\arrow[l,"{\substack{\text{induit}\\\text{par }p|{\widetilde f}:\ \widetilde f\simeq f}}"]
\arrow[r,"{\alpha_{\widetilde f}}"] \& (\det M)^*
\end{tikzcd}}
est indépendante de $\widetilde f$, soit
\[
\alpha_f:\ \underline{\omega}_f\longrightarrow(\det M)^*\tag*{(31)}
\]
et on a
\[
SF_M=\{f\in F_M\mid\alpha_f(\underline{\omega}_f)\unsure{\subset|=}\{\pm e_M\}\}\tag*{(32)}
\]
Soit maintenant
\[
a\in\deleted{S}A_M\qquad a=\{\dot x,\dot y\},\ \text{avec }(x,y)\in\mathrm{Rep}\,M
\]
et soient $f,f'$ les deux éléments de $F_M$ contenant $a$, donc
\[
\begin{aligned}
f&=\{\dot x,\dot y,(-x-y)^{\cdot}\}\\
f'&=\{\dot x',\deleted{\unsure{\dot y}}\ \dot y=(-y)^{\cdot},(-x+y)^{\cdot}\}
\end{aligned}
\]
Considérons les deux applications
\diagram{\begin{tikzcd}[ampersand replacement=\&]
\deleted{a}\arrow[r,"\deleted{\unsure{\sim}}"] \& \deleted{\underline{\omega}_f}\arrow[r,"\deleted{\alpha_f}"] \& \deleted{(\det M)^*}
\end{tikzcd}}

\diagram{\[\begin{tikzcd}[ampersand replacement=\&,row sep=small]
\& \underline{\omega}_f\arrow[dr,"\alpha_f"] \& \\
a\arrow[ur,"{\substack{\text{can}\\\sim}}"]\arrow[dr,no head,"\sim","{\text{can}}"'] \& \&(\det M)^*\quad\unsure{)|,}\\
\&\underline{\omega}_{f'}\arrow[ur,"\alpha_{f'}"]\&
\end{tikzcd}\tag*{(33)}\]}
% archive sheet number: 12
